\section*{Schedules}
\label{sec:schedules}

\scheduleSection{Groups}{8}{Basic concepts of group theory recalled from Part IA Groups. Normal subgroups, quotient groups and isomorphism theorems. Permutation groups. Groups acting on sets, permutation representations. Conjugacy classes, centralizers and normalizers. The centre of a group. Elementary properties of finite \(p\)-groups. Examples of finite linear groups and groups arising from geometry. Simplicity of \(A_n\).

\noindent Sylow subgroups and Sylow theorems. Applications, groups of small order.}

\scheduleSection{Rings}{10}{Definition and examples of rings (commutative, with 1). Ideals, homomorphisms, quotient rings, isomorphism theorems. Prime and maximal ideals. Fields. The characteristic of a field. Field of fractions of an integral domain.

\noindent Factorization in rings; units, primes and irreducibles. Unique factorization in principal ideal domains, and in polynomial rings. Gauss' Lemma and Eisenstein's irreducibility criterion.

\noindent Rings \(\ZZ\brac{\alpha}\) of algebraic integers as subsets of \(\CC\) and quotients of \(\ZZ\brac{x}\). Examples of Euclidean domains and uniqueness and non-uniqueness of factorization. Factorization in the ring of Gaussian integers; representation of integers as sums of two squares.

\noindent Ideals in polynomial rings. Hilbert basis theorem.}

\scheduleSection{Modules}{6}{Definitions, examples of vector spaces, abelian groups and vector spaces with an endomorphism. Submodules, homomorphisms, quotient modules and direct sums. Equivalence of matrices, canonical form. Structure of finitely generated modules over Euclidean domains, applications to abelian groups and Jordan normal form.}